\section{Component Analysis of Two-Phase Hybrids}
\label{sec:ablation}
% SWEVO extension: two generated diagnostics tables (tab:D-vns, tab:D-rsreplay) are placed in this section; they are
% captured from analysis/mpce_supp_diag.tex, whose figure (fig:D-psodyn) is discarded here (it is not placed by this section)
\begingroup\RenewEnviron{figure*}[1][]{}\CaptureSuppTables{analysis/mpce_supp_diag.tex}\endgroup
% posterior probabilities of the Bayesian signed-rank test in math mode: "= value", or "< 0.01" / "> 0.99" for a bound
% (same rule as \BayEq of the supplement; own name, so that no other section file can clash with it)
\makeatletter
\newcommand{\AblBayEq}[1]{\expandafter\abl@bayeq#1\@nil}
\def\abl@bayeq#1#2\@nil{\ifx#1\ensuremath #2\else =#1#2\fi}
\makeatother
%% generated by mpce_results.py -- do not edit by hand
\begin{table}[!t]
\centering
\caption{Ablation over the 68 Benchmark Cases (30 Seed-Paired Runs, 6,030 Calls). Top: Average Rank Among the Eight Variants (Friedman $\chi^2_F=326.8$, 7 d.f., $p=1.14\times10^{-66}$) and Feasible Runs. Bottom: Planned Contrasts, Cases in Which the First Variant Is Significantly Better / Not Different / Worse (Wilcoxon, Holm-Adjusted over the Eight Contrasts of Each Case, $\alpha=0.05$), and Mean Difference of the Wake Loss $\overline{\Delta L}$ (Percentage Points; Negative = First Variant Better)}
\label{tab:ablation}
\scriptsize\setlength{\tabcolsep}{2.5pt}
\begin{tabular}{lcccc}
\toprule
Variant & Phase 1 & Phase 2 & Avg.\ rank & Feas.\ (\%) \\
\midrule
PSO-VNS & PSO & VNS & 1.87 & 100.0 \\
SSA-VNS & SSA & VNS & 3.79 & 99.6 \\
LX-SSA-VNS & LX-SSA & VNS & 4.62 & 99.9 \\
RS-VNS & random sampling & VNS & 4.46 & 99.4 \\
VNS & -- & VNS & 5.32 & 100.0 \\
PSO & PSO & -- & 2.10 & 100.0 \\
SSA & SSA & -- & 6.56 & 97.6 \\
LX-SSA & LX-SSA & -- & 7.29 & 96.9 \\
\bottomrule
\end{tabular}

\smallskip
\begin{tabular}{l>{\raggedright\arraybackslash}p{2.9cm}cc}
\toprule
Contrast & Isolates & W/T/L & $\overline{\Delta L}$ \\
\midrule
PSO-VNS vs.\ PSO & VNS phase & 8/58/2 & $-0.018$ \\
PSO-VNS vs.\ VNS & swarm start vs.\ best initial point & 50/17/1 & $-0.465$ \\
PSO-VNS vs.\ RS-VNS & swarm vs.\ random sampling & 48/20/0 & $-0.399$ \\
PSO-VNS vs.\ SSA-VNS & PSO vs.\ SSA as Phase~1 & 47/21/0 & $-0.330$ \\
PSO-VNS vs.\ LX-SSA-VNS & PSO vs.\ LX-SSA as Phase~1 & 48/20/0 & $-0.393$ \\
SSA-VNS vs.\ RS-VNS & SSA vs.\ random sampling & 2/66/0 & $-0.068$ \\
SSA-VNS vs.\ LX-SSA-VNS & Laplace step in the hybrid & 3/65/0 & $-0.063$ \\
LX-SSA vs.\ SSA & Laplace step alone & 0/61/7 & $+0.166$ \\
\bottomrule
\end{tabular}
\end{table}


The variants of Table~\ref{tab:ablation} combine a Phase~1 (PSO, SSA, LX-SSA, random sampling or none) with or without the same VNS phase (6,030 evaluations from random starts, same seeds; convergence: Fig.~\ref{fig:ablation}); their average ranks differ (Friedman $p=\NAblP$). W/T/L, $p$ (Holm-adjusted Wilcoxon test on the 68 case means) and $\overline{\Delta L}$ (mean difference of the case means, pp) are those of Table~\ref{tab:ablation}; a 90\% bootstrap interval of $\overline{\Delta L}$ (in brackets) inside $\pm\NEqMargin$~pp means practical equivalence at the case level.
% CHECK-FINAL [C36]: Friedman rejects equal average ranks (\NAblP); pipeline-generated counts \NAblNContrasts (Holm family of each case), \NAblDf d.f., \NAblNVar variants in the caption. All DL in this section are \NCm...DL = tab:ablation values (R1-7, R5-F1)
Each contrast changes one component and keeps the other fixed, on seed-paired runs that share the initial population (Section~\ref{sec:setup}), and so answers one question. A Phase~1 with and without the VNS phase measures the value of the local-search phase (Section~\ref{sec:abl-vns}). A Phase~1 against a random-sampling control of equal cost, both followed by the same VNS phase, measures whether the Phase-1 method searches better than drawing layouts at random (Sections~\ref{sec:abl-controls} and~\ref{sec:abl-swarm}). Two Phase-1 methods followed by the same VNS phase show which of them gives VNS the better start (Section~\ref{sec:abl-swarm}). Section~\ref{sec:abl-equiv} asks which differences are of practical size, and Section~\ref{sec:split} varies the share of the budget given to Phase~1. Table~\ref{S-tab:ablation-cases} gives the per-case values and outcomes of the \NAblNContrasts{} contrasts.
% SWEVO extension (design of the contrasts; no outcome stated). CHECK-FINAL [C36]: \NAblNContrasts contrasts per case (pipeline count)

\subsection{The VNS Phase}
\label{sec:abl-vns}
The VNS phase improves both salp-swarm methods: SSA-VNS is \NAblSSAVNSvsSSA{} (W/T/L) against SSA and LX-SSA-VNS \NAblLXSSAVNSvsLXSSA{} against LX-SSA ($\overline{\Delta L}=\NCmSSAVNSvsSSADL$ and \NCmLXSSAVNSvsLXSSADL). For PSO, which PSO-VNS repeats for its first 3,030 evaluations (Section~\ref{sec:hybrid}), the difference is small, not significant and practically equivalent to zero ($\overline{\Delta L}=\NCmPSOVNSvsPSODL$ \NEqPSOVNSvsPSOCI, $p=\NCmPSOVNSvsPSOPHolm$). At this budget the VNS phase is essentially one compass-search descent: for PSO-VNS, the shake-and-search cycles take \NDCycleEvalsPctPSOVNS\% of the Phase-2 evaluations (pooled; median run \NDCycleEvalsPctMedianPSOVNS\%) but give \NDCycleSharePSOVNS\% of the Phase-2 gain (\NDCycleAllSharePSOVNS\% with the cycle cut off by the budget), and the median run with $N\ge12$ never shakes (Table~\ref{tab:D-vns}). The analysis thus measures the value of a local-search phase.
% CHECK-FINAL [C33]: VNS phase significant for SSA and LX-SSA (W > L, L = 0; case-mean p < 0.05, DL < 0), not for PSO (case-mean p >= 0.05, so also Holm); [C13]: PSOBV-PSOC W > L; [C49]: case-level 90 % CI \NEqPSOVNSvsPSOCI inside +-\NEqMargin
% CHECK-DIAG [D20]: cycles >= 15 % of the Phase-2 evaluations, < 5 % of the gain incl. the cut-off cycle, accepted cycles < 1 % (\NDCycleEvalsPctPSOVNS, median run \NDCycleEvalsPctMedianPSOVNS, \NDCycleSharePSOVNS, \NDCycleAllSharePSOVNS); [D08]: essentially one descent; [D09]: median run has no shaking step for N >= 12

\SuppTable{tab:D-vns}
Table~\ref{tab:D-vns} follows the VNS phase of PSO-VNS, SSA-VNS and RS-VNS on the \NDVnsCases{} cases of the budget-split study (Section~\ref{sec:split}) $\times$ \NDVnsSeeds{} seeds. It separates the \emph{first descent}, the compass search of Section~\ref{sec:vns} from the switch point until the step size falls below its lower limit, from the \emph{cycles}, each a shaking step followed by a complete local search. In PSO-VNS, the first descent takes a median of \NDDescentEvalsPSOVNS{} evaluations (pooled share \NDDescentEvalsPctPSOVNS\% of the Phase-2 evaluations, median share per run \NDDescentEvalsPctMedianPSOVNS\%) and yields \NDDescentSharePSOVNS\% of the Phase-2 improvement; the median number of shaking steps per run is \NDShakesPSOVNS{} (range \NDShakesMinPSOVNS--\NDShakesMaxPSOVNS), and \NDNoAcceptedPSOVNS{} of the \NDVnsRunsPSOVNS{} runs have no accepted cycle. SSA-VNS and RS-VNS follow the same pattern (first descent \NDDescentShareSSAVNS\% and \NDDescentShareRSVNS\% of the improvement; \NDNoAcceptedSSAVNS{} and \NDNoAcceptedRSVNS{} runs without an accepted cycle).
% CHECK-DIAG [D08]: for PSO-VNS, SSA-VNS and RS-VNS: median shaking steps <= 2, shaking < 1 % of the Phase-2 evaluations, >= 85 % of the Phase-2 improvement from the first descent (\NDDescentSharePSOVNS, \NDDescentShareSSAVNS, \NDDescentShareRSVNS); no accepted cycle in \NDNoAcceptedPSOVNS / \NDNoAcceptedSSAVNS / \NDNoAcceptedRSVNS of 120 runs
The cycles are expensive although a shake costs a single evaluation: each shake is followed by a new compass search over all $2n$ coordinate moves, so the cycles take \NDCycleEvalsPctImpPSOVNS\% of the Phase-2 evaluations of the PSO-VNS runs with a feasible switch point and return little of the gain. The first descent needs more evaluations as $N$ grows, since each compass iteration costs $2n=4N$ evaluations: for $N\ge12$ it is unfinished within the budget in \NDDescentUnfinishedLargePSOVNS{} of the \NDRunsLargePSOVNS{} PSO-VNS runs, and the median run never reaches a shaking step (median number of shaking steps \NDShakesNSixPSOVNS{} for $N=6$, \NDShakesNFifteenPSOVNS{} for $N=15$). All accepted cycles use the smallest neighborhood $k=1$; $k_{\max}=5$ is never reached.
% CHECK-DIAG [D20]: cycles take more than twice their gain share in the runs with a feasible switch point (\NDCycleEvalsPctImpPSOVNS); [D09]: N >= 12: first descent unfinished in >= 90 % of the runs (\NDDescentUnfinishedLargePSOVNS of \NDRunsLargePSOVNS), median run without shaking step; [D10]: accepted cycles only with k = 1, k_max never reached
The VNS phase never makes a feasible incumbent infeasible; it repairs infeasible switch points only in the first descent, and its start is always the best feasible Phase-1 layout when Phase~1 evaluated one, so a control that hands VNS the best random sample would coincide with RS-VNS.
% CHECK-DIAG [D11]: VNS never turns a feasible incumbent infeasible; infeasible switch points repaired only by the first descent; [D12]: the start is the best feasible Phase-1 layout ("rsbest" control = RS-VNS)
At 6,030 evaluations the component analysis therefore measures the value of a truncated best-improvement compass search after Phase~1. The shaking and the neighborhood structure of VNS can matter only for small $N$, and for $N\ge12$ the local-minimality guarantee of the compass search (Proposition~\ref{S-prop:S-ls}) does not apply, as its descent is cut off by the budget.
% CHECK-DIAG [D08]-[D10]: interpretation of the diagnostics above (one descent; shaking only for small N; unfinished descent for N >= 12)

\subsection{Random-Sampling Controls}
\label{sec:abl-controls}
Random sampling of equal cost replaces Phase~1, uniform in the site disc (RSD-VNS) or in its bounding square (RS-VNS); VNS alone is a third control. RS-VNS is weak: \NRSRunsNoFeasSample{} of its runs draw no feasible sample, although all of them draw samples wholly inside the circle: the spacing constraint is binding, and disc sampling rescues only \NRsSquareExplains{} of these runs. RSD-VNS beats RS-VNS (\NAblRSDVNSvsRSVNS, $\overline{\Delta L}=\NCmRSDVNSvsRSVNSDL$) and ranks third (\NRankRSDVNS).
% CHECK-FINAL [C60]: RS-VNS runs without a feasible sample (\NRSRunsNoFeasSample) = rescued by disc sampling (\NRsSquareExplains) + none either way (\NRsSpacingDominates), rescued < none; LEAD: extend C60 with runs_no_feasible_but_inside_sample == runs_no_feasible_sample (\NRsNoFeasInside = \NRSRunsNoFeasSample, mpce_results.py phase1_replay.RSVNS) for "all of them"
% CHECK-FINAL [C57]: RSDVNS-RSVNS L == 0, W > 0, case-mean p_Holm < 0.05, DL < 0; [C58]: RSD-VNS ranks third of \NAblNVar

\SuppTable{tab:D-rsreplay}
Why RS-VNS is weak follows from Table~\ref{tab:D-rsreplay}, which replays Phase~1 of RS-VNS with the seeded random streams of the runs (geometry only; the replay agrees with the stored runs). A square sample has all $N$ turbines inside the circle with probability $(\pi/4)^N$ (Proposition~\ref{S-prop:S-box}), and the observed inside rates match it; yet the feasible samples vanish much sooner. For $N=15$ in the 1000-m farm, \NDRsInsideNFifteen\% of the \NDRsSamples{} samples lie wholly inside the circle ($(\pi/4)^{15}$: \NDRsInsideTheoryNFifteen\%), but none is feasible, because random positions rarely keep all $N(N-1)/2$ pairs at least $\ell_{\min}$ apart.
% CHECK-DIAG [D17]: replay agrees with the runs; inside rate matches (pi/4)^N; N = 15, 1000 m: \NDRsFeasNFifteen of \NDRsSamples box samples feasible (inside \NDRsInsideNFifteen vs theory \NDRsInsideTheoryNFifteen)
In \NRSRunsNoFeasSample{} RS-VNS runs, among them all runs of \NDRsCasesNoFeasSample{} cases, Phase~1 draws no feasible layout at all, and VNS starts from the least-penalized sample (Section~\ref{S-sec:S-th-budget}). Every one of these runs does draw samples with all turbines inside the circle (\NRsNoFeasInside{} of \NRSRunsNoFeasSample), so the spacing constraint, not the boundary, is binding in all of them. Paired by seed with the disc control, \NRsSquareExplains{} of these runs obtain a feasible sample in the disc, whereas \NRsSpacingDominates{} have none under either sampling. The weakness of RS-VNS thus lies mainly in the spacing constraint and only in part in the square, and a Phase-1 method that beats RS-VNS passes only a minimum bar. The last column of Table~\ref{tab:D-rsreplay} gives the feasibility rate of samples drawn uniformly in the disc, the Phase-1 distribution of RSD-VNS.
% CHECK-DIAG [D17]: \NDRsRunsNoFeasSample runs (all runs of \NDRsCasesNoFeasSample cases) without a feasible Phase-1 sample; CHECK-FINAL [C60]: \NRsNoFeasInside = \NRSRunsNoFeasSample ("every one"), rescued \NRsSquareExplains < none either way \NRsSpacingDominates ("mainly in the spacing")

\SuppTable{tab:switch}
Table~\ref{tab:switch} describes the hand-over from Phase~1 to VNS in all two-phase variants: when the switch occurs, how often the layout handed to VNS is feasible, its wake loss, which share of the remaining wake loss the VNS phase removes, and how many runs infeasible at the switch end feasible. The controls switch after 3,015 evaluations, the swarm hybrids after $N_p(T_1+1)=3{,}030$ (Section~\ref{sec:template}); for the controls, feasibility at the switch is read at the last stored checkpoint (3,000 evaluations), which for RS-VNS gives the same runs as at 3,015. The mean loss at the switch is taken over the runs feasible at that point, so it is comparable only between variants that are feasible in about the same runs.
% CHECK-FINAL [C61]: controls switch at round(0.5 x 6,030) = 3,015, swarms at 3,030; RS-VNS runs infeasible at 3,000 = replayed runs without a feasible sample among all 3,015
Sampling in the disc raises the share of feasible Phase-1 samples from \NRSPctFeasSamples\% to \NRSDPctFeasSamples\% and lowers the number of runs without any feasible Phase-1 sample from \NRSRunsNoFeasSample{} to \NRSDRunsNoFeasSample; at the switch, the best layout is feasible in \NSwitchFeasRSDVNS\% of the RSD-VNS runs against \NSwitchFeasRSVNS\% of the RS-VNS runs. Both controls remain well below the swarm phases (PSO \NSwitchFeasPSOVNS\%, SSA \NSwitchFeasSSAVNS\%), so their VNS phase starts more often from an infeasible layout and has to repair it first.
% CHECK-FINAL [C57]: RSD-VNS Phase 1 feasible at the switch more often and fewer runs without a feasible sample (\NRSDRunsNoFeasSample < \NRSRunsNoFeasSample; replay verified); CHECK-DIAG [D18]: feasible by the switch PSO > SSA > RS (RSD-VNS \NSwitchFeasRSDVNS < SSA-VNS \NSwitchFeasSSAVNS: tab:switch values); [D11]: infeasible switch points repaired by the first descent

\subsection{Swarm First Phases}
\label{sec:abl-swarm}
A PSO phase clearly beats every control: PSO-VNS is \NAblPSOVNSvsRSDVNS{} against RSD-VNS ($\overline{\Delta L}=\NCmPSOVNSvsRSDVNSDL$ \NEqPSOVNSvsRSDVNSCI, wholly below $-\NEqMargin$), \NAblPSOVNSvsRSVNS{} against RS-VNS and \NAblPSOVNSvsVNS{} against VNS. At this budget and from random starts, salp-swarm phases add nothing beyond random sampling in the farm. Against RSD-VNS, SSA-VNS is \NAblSSAVNSvsRSDVNS{} ($\overline{\Delta L}=\NCmSSAVNSvsRSDVNSDL$; Hodges--Lehmann estimate \NXHLSSAVNSvsRSDVNS~pp, 95\% CI \NXHLSSAVNSvsRSDVNSCI; the case means, not the farm clusters, favor RSD-VNS): a gain of the SSA phase larger than the margin is ruled out at the seed and the case level (90\% intervals \NXEqSeedSSAVNSvsRSDVNSCI{} and \NXEqCaseSSAVNSvsRSDVNSCI, both above $-\NEqMargin$), whereas equivalence remains inconclusive. LX-SSA-VNS is \NAblLXSSAVNSvsRSDVNS{} and worse ($\overline{\Delta L}=\NCmLXSSAVNSvsRSDVNSDL$, $p=\NCmLXSSAVNSvsRSDVNSPHolm$). The small SSA gain over RS-VNS (\NAblSSAVNSvsRSVNS, $\overline{\Delta L}=\NCmSSAVNSvsRSVNSDL$) only reflects that control's weakness; LX-SSA-VNS and RS-VNS are practically equivalent (\NEqLXSSAVNSvsRSVNSCI). Re-evaluated with 1$^\circ$ direction bins (Table~\ref{S-tab:F-abl}), SSA-VNS and LX-SSA-VNS again have a higher loss than RSD-VNS (\NFAblSSAVNSvsRSDVNSOne{} and \NFAblLXSSAVNSvsRSDVNSOne~pp) and PSO-VNS beats all three controls; only the SSA gain over RS-VNS is no longer significant.
% CHECK-FINAL [C56]: PSOBV-RSDVNS W >= 40, L == 0, p_Holm < 0.05, whole 90 % CI below -\NEqMargin; [C14]: PSOBV-RSVNS L == 0, PSOBV-BVNS L <= 2 ("clearly beats every control")
% CHECK-FINAL [C53]: SSABV-RSDVNS W == 0, case-mean losses > wins, p_Holm < 0.05 against SSA-VNS; CHECK-EXTRA [X33]: not at cluster level ("not the farm clusters"; never "worse"); [X49]: HL > 0, 95 % CI \NXHLSSAVNSvsRSDVNSCI; [X53]: lower 90 % limits above -m at seed and case level, not at cluster level ("ruled out at the seed and the case level"); CHECK-FINAL [C54]: 90 % case CI not inside +-m (equivalence not shown); CHECK-EXTRA [X54]: seed level equivalent, case/cluster not ("inconclusive")
% CHECK-FINAL [C55]: LXBV-RSDVNS L > W, p_Holm < 0.05, DL > 0; CHECK-EXTRA [X34]: CR p 0.087 ("worse", not "clearly worse"); CHECK-FINAL [C15]/[C51]: SSABV-RSVNS small gain, not equivalent; [C50]: LX-SSA-VNS vs RS-VNS equivalent (\NEqLXSSAVNSvsRSVNSCI)
% CHECK-DIR [F16]: 1-deg: SSABV and LXBV higher case-mean loss than RSDVNS (p < 0.05), PSOBV better than RSDVNS, RSVNS, BVNS (p < 0.05); [F17]: 1-deg SSABV-RSVNS p >= 0.05 (\NFAblSSAVNSvsRSVNSOneP)
The case means and the farm clusters answer different questions (Section~\ref{sec:setup}): the case-level tests treat the 68 cases as exchangeable, whereas the six clusters (data set and farm radius, with nested $N$) allow a generalization to new farms. The disadvantage of SSA-VNS against RSD-VNS holds at the case level (Holm-adjusted $p=\NCmSSAVNSvsRSDVNSPHolm$) but not over the clusters (RSD-VNS lower in \NXClSSAVNSvsRSDVNSOpp{} of \NXClustersNum, cluster-robust $p=\NXClSSAVNSvsRSDVNSCRP$; Section~\ref{S-sec:S-inference}); we therefore do not call SSA-VNS worse than disc sampling, only not better.
% CHECK-FINAL [C53]: case-mean p_Holm < 0.05 against SSA-VNS; CHECK-EXTRA [X33]: does NOT survive the cluster analyses (RSD-VNS lower in \NXClSSAVNSvsRSDVNSOpp of six clusters, CR p \NXClSSAVNSvsRSDVNSCRP) -> "not better", never "worse"

\emph{LX-SSA as published.} LX-SSA costs $2N_p$ evaluations per iteration (Section~\ref{sec:ssa}) and so runs half as many iterations as SSA; the contrasts do not separate its Laplace step from this cost. As published, it makes the salp swarm worse: LX-SSA is \NAblLXSSAvsSSA{} against SSA ($\overline{\Delta L}=\NCmLXSSAvsSSADL$, $p=\NCmLXSSAvsSSAPHolm$) and SSA-VNS is \NAblSSAVNSvsLXSSAVNS{} against LX-SSA-VNS ($\overline{\Delta L}=\NCmSSAVNSvsLXSSAVNSDL$, $p=\NCmSSAVNSvsLXSSAVNSPHolm$).
% CHECK-FINAL [C17]: LX-SSA as published worse (never "Laplace step harmful"): LXSSA-SSA W == 0; SSABV-LXBV L == 0; case means favour SSA and SSA-VNS (p < 0.05, so also Holm)
As published, the Laplace candidate thus does not make up for the iterations it costs; whether the Laplace step would help at equal iterations is not tested here. The VNS phase narrows but does not close the gap: after it, LX-SSA-VNS is still worse than SSA-VNS, and it is no better than the square-sampling control (Section~\ref{sec:abl-equiv}).
% CHECK-FINAL [C17]: SSABV-LXBV L == 0, case means favour SSA-VNS ("still worse"); [C35]: LX-SSA-VNS vs RS-VNS case-mean p >= 0.05 ("no better"); [C50]: equivalent. "Narrows": |DL| SSA-VNS vs LX-SSA-VNS (\NCmSSAVNSvsLXSSAVNSDL) < |DL| LX-SSA vs SSA (\NCmLXSSAvsSSADL), tab:ablation values; CHECK-EXTRA [X38]: tab:equivalence rows reproduced

\emph{Choice of Phase~1.} With the same VNS phase, PSO-VNS is \NAblPSOVNSvsSSAVNS{} against SSA-VNS and \NAblPSOVNSvsLXSSAVNS{} against LX-SSA-VNS. At the switch, its global best is feasible more often (\NSwitchFeasPSOVNS\% of the runs vs.\ \NSwitchFeasSSAVNS\% and \NSwitchFeasLXSSAVNS\%) and, when feasible, has a lower wake loss (\NSwitchLossPSOVNS\% vs.\ \NSwitchLossSSAVNS\% and \NSwitchLossLXSSAVNS\%). Whether PSO helps mainly by faster feasibility or by better exploration remains open (cf.\ the feasible starts in Section~\ref{sec:feasinit}).
% CHECK-FINAL [C16]: PSOBV-SSABV, PSOBV-LXBV: L == 0, W >= 34; switch feasibility higher and switch loss lower for PSOBV
As VNS is essentially one descent from the switch point (Section~\ref{sec:abl-vns}), the start it receives largely decides the final layout, and the three swarm hybrids differ only in that start.
% CHECK-DIAG [D08]: first descent >= 85 % of the Phase-2 improvement for PSO-VNS and SSA-VNS (interpretation of tab:D-vns)

\subsection{Practical Equivalence of the Contrasts}
\label{sec:abl-equiv}
A non-significant test does not show that two variants are equally good, and a significant one does not show that the difference matters. Table~\ref{tab:equivalence} therefore gives, for every contrast, the size of the difference of the case means with its 95\% and 90\% bootstrap intervals, the TOST decision~\cite{Lakens2017} at the margin $m=\NEqMargin$~pp, the minimal margin $m_{\min}$ at which equivalence would hold, and the posterior probabilities of the Bayesian signed-rank test~\cite{Benavoli2017}: the first variant better by more than $m$ ($P_{\rm A}$), the two within $\pm m$ ($P_{\rm rope}$), the second better ($P_{\rm B}$). The margin is a convention set after the primary analysis (Section~\ref{sec:setup}); $m_{\min}$ lets readers apply their own. A pair can be equivalent (the whole 90\% interval inside $\pm m$), different by more than the margin (the whole interval beyond $-m$ or $+m$), or neither, when the interval crosses $-m$ or $+m$; the component contrasts show all three outcomes.
% CHECK-FINAL [C49]-[C52]: TOST = 90 % bootstrap CI inside (-m, m); m = \NEqMargin; Bayesian signed-rank with rope (-m, m); CHECK-EXTRA [X38]: all 18 pairs of tab:equivalence (CIs, TOST, m_min, verdict, Bayes) reproduced from the per-run data; [X56]: margin a post hoc convention

\SuppTable{tab:equivalence}
\emph{Equivalent.} Two contrasts are practically equivalent at the case level. The VNS phase after PSO changes the case means by less than the margin (90\% CI \NEqPSOVNSvsPSOCI~pp, $m_{\min}=\NEqPSOVNSvsPSOMin$~pp, $P_{\rm rope}\AblBayEq{\NBayPSOVNSvsPSORope}$); Section~\ref{sec:overall} discusses this pair and the level of inference at which it holds. LX-SSA-VNS and RS-VNS are equivalent as well (\NEqLXSSAVNSvsRSVNSCI~pp, $m_{\min}=\NEqLXSSAVNSvsRSVNSMin$~pp, $P_{\rm rope}\AblBayEq{\NBayLXSSAVNSvsRSVNSRope}$), but not over the farm clusters ($m_{\min}=\NXEqClustLXSSAVNSvsRSVNSMin$~pp); no pair of Table~\ref{tab:equivalence} is equivalent at the cluster level.
% CHECK-FINAL [C49]: PSO-VNS vs PSO equivalent at the case level, P(rope) > 0.99; [C50]: LX-SSA-VNS vs RS-VNS equivalent; CHECK-EXTRA [X55]: LX-SSA-VNS vs RS-VNS not equivalent at the cluster level (\NXEqClustLXSSAVNSvsRSVNSMin); no pair of tab:equivalence equivalent at the cluster level; [X39]: PSO-VNS vs PSO not equivalent at the cluster level (discussed in sec:overall)

\emph{Different by more than the margin.} PSO-VNS gains more than the margin over every control: its whole 90\% interval lies below $-\NEqMargin$~pp against RSD-VNS (\NEqPSOVNSvsRSDVNSCI~pp, $P_{\rm A}\AblBayEq{\NBayPSOVNSvsRSDVNSBetter}$) and against RS-VNS (\NEqPSOVNSvsRSVNSCI~pp). The VNS phase after SSA and after LX-SSA, and PSO-VNS against VNS and against the other methods of the main comparison (lower block of the table), also give differences whose whole 90\% interval lies beyond the margin. So does LX-SSA as published against SSA, with its whole interval above $+\NEqMargin$~pp: without the VNS phase, LX-SSA is worse than SSA by more than the margin. Among the controls, the disc control is better than the square control (RSD-VNS $-$ RS-VNS: \NEqRSDVNSvsRSVNSCI~pp, $P_{\rm A}\AblBayEq{\NBayRSDVNSvsRSVNSBetter}$).
% CHECK-FINAL [C56]: PSO-VNS vs RSD-VNS whole 90 % CI below -m; [C14]: PSO-VNS beats RS-VNS and VNS (CIs \NEqPSOVNSvsRSVNSCI and [-0.545, -0.386] below -m, tab:equivalence); [C33]: VNS phase significant for SSA and LX-SSA (90 % CIs [-0.368, -0.226] and [-0.505, -0.293] below -m in tab:equivalence); [C17]: LX-SSA as published worse (90 % CI of LX-SSA vs SSA [0.104, 0.206] above +m in tab:equivalence); [C57]: RSD-VNS better than RS-VNS (90 % CI \NEqRSDVNSvsRSVNSCI, upper limit below -m); CHECK-EXTRA [X38]: tab:equivalence CIs, verdicts and Bayes probabilities reproduced (main block: all five PSO-VNS rows below -m)

\emph{Neither.} The remaining component contrasts are neither equivalent nor shown to differ by more than the margin. The small gain of SSA-VNS over the weak control RS-VNS is not equivalent to zero (\NEqSSAVNSvsRSVNSCI~pp, $m_{\min}=\NEqSSAVNSvsRSVNSMin$~pp), but a negligible gain is the more probable outcome in the Bayesian analysis ($P_{\rm rope}\AblBayEq{\NBaySSAVNSvsRSVNSRope}$ against $P_{\rm A}\AblBayEq{\NBaySSAVNSvsRSVNSBetter}$). After the VNS phase, the disadvantage of LX-SSA shrinks: for SSA-VNS against LX-SSA-VNS (\NEqSSAVNSvsLXSSAVNSCI~pp, $m_{\min}=\NEqSSAVNSvsLXSSAVNSMin$~pp), a gain of practical size and a negligible difference are about equally probable ($P_{\rm A}\AblBayEq{\NBaySSAVNSvsLXSSAVNSBetter}$, $P_{\rm rope}\AblBayEq{\NBaySSAVNSvsLXSSAVNSRope}$), and LX-SSA-VNS is significantly better in no case.
% CHECK-FINAL [C51]: SSA-VNS vs RS-VNS not equivalent, P(rope) \NBaySSAVNSvsRSVNSRope > P(better) \NBaySSAVNSvsRSVNSBetter; [C17]: SSABV-LXBV L == 0 ("significantly better in no case"), case means favour SSA-VNS; CHECK-EXTRA [X38]: SSA-VNS vs LX-SSA-VNS not equivalent, 90 % CI \NEqSSAVNSvsLXSSAVNSCI crosses -m, P_A \NBaySSAVNSvsLXSSAVNSBetter vs P_rope \NBaySSAVNSvsLXSSAVNSRope (tab:equivalence reproduced)
For SSA-VNS against RSD-VNS, the question that matters is one-sided: does the SSA phase give VNS a better start than disc sampling by a practically relevant amount? A gain of the SSA phase larger than the margin is ruled out at the seed and the case level, where the lower limits of the 90\% intervals of SSA-VNS $-$ RSD-VNS lie above $-\NEqMargin$~pp (\NXEqSeedSSAVNSvsRSDVNSCI{} and \NXEqCaseSSAVNSvsRSDVNSCI~pp), but not at the cluster level (CR2 interval \NXEqClustSSAVNSvsRSDVNSCI~pp). Equivalence itself is not shown at the case level ($m_{\min}=\NEqSSAVNSvsRSDVNSMin$~pp), and the difference points against SSA-VNS ($P_{\rm A}\AblBayEq{\NBaySSAVNSvsRSDVNSBetter}$, $P_{\rm rope}\AblBayEq{\NBaySSAVNSvsRSDVNSRope}$, $P_{\rm B}\AblBayEq{\NBaySSAVNSvsRSDVNSWorse}$). We therefore state this result one-sidedly: at this budget and from random starts, an SSA phase gives no gain of practical size over sampling in the disc. LX-SSA-VNS is not equivalent to RSD-VNS and worse than it (90\% CI \NEqLXSSAVNSvsRSDVNSCI~pp, $P_{\rm B}\AblBayEq{\NBayLXSSAVNSvsRSDVNSWorse}$), with an interval that lies above zero but reaches below $+\NEqMargin$~pp.
% CHECK-EXTRA [X53]: lower 90 % limits of SSA-VNS - RSD-VNS above -m at the seed and case levels, not at the cluster level (CR2 \NXEqClustSSAVNSvsRSDVNSCI); CHECK-FINAL [C54]: equivalence not shown at the case level, P(SSA-VNS better) < 0.01, P(rope) > 0.5; [C55]: LX-SSA-VNS worse than RSD-VNS, 95 % CI above 0, not equivalent; CI \NEqLXSSAVNSvsRSDVNSCI: 0 < lower limit < m ("reaches below +m")

\subsection{Budget Split}
\label{sec:split}
\CaptureSuppTables{analysis/mpce_tab_split.tex}
\SuppTable{tab:split}
The split $\omega$ is the share of the 6,030 evaluations given to PSO; the rest goes to VNS. The study compares the preset $\omega=0.5$, fixed before the experiments and used for all other results, with $\omega=0.25$, $0.75$ and $0.9$ and with the end point $\omega=1$, plain PSO on the same seeds (\NSplitNSettings{} settings), on the \NSplitCases{} cases with the middle and the largest $N$ of each farm and data set (Table~\ref{tab:split}; per case: Table~\ref{S-tab:split-cases}). With $\omega=0.75$, the mean wake loss is lower than with $\omega=0.5$ in \NSplitSeventyFiveLower{} cases (significantly in \NSplitSeventyFiveSig; by \NSplitGainSeventyFive~pp on average) and significantly higher in \NSplitFiftyBetterSeventyFive{} case; $\omega=0.25$ is significantly worse than the preset in \NSplitTwentyFiveSig{} cases and better in \NSplitTwentyFiveBetter{} case. The loss is thus lowest at an interior split: shortening the PSO phase and dropping the VNS phase both raise it. Plain PSO has a higher mean loss than $\omega=0.75$ (\NSplitLossHundred\% vs.\ \NSplitLossSeventyFive\%, average rank \NSplitRankHundred{} vs.\ \NSplitRankSeventyFive; W/T/L of $\omega=0.75$: \NSplitHundredVsSeventyFive), so the best tested split is not the PSO end point. These conclusions hold in all six farm clusters.
% CHECK-FINAL [C18]: 75 % lower loss than 50 % in >= 10 of 12 cases, 50 % never significantly better than 75 %, 25 % never significantly better than 50 %, ranks 75 < 50 < 25; [C47]: 0.75 lowest mean loss and rank of the \NSplitNSettings settings, omega = 1 higher mean loss than 0.75 ("not the PSO end point", "interior split"); CHECK-EXTRA [X37]: 0.5 better than 0.25, 0.75 better than 0.5, 0.75 better than 1 hold in all six clusters (0.5 vs 1 splits the clusters: not stated)
Moving further toward PSO does not help: $\omega=0.9$ differs significantly from $\omega=0.75$ in no case (W/T/L \NSplitNinetyVsSeventyFive, case means $p=\NSplitCmNinetyP$) and has a slightly higher mean loss (\NSplitLossNinety\% vs.\ \NSplitLossSeventyFive\%) and a worse average rank (\NSplitRankNinety{} vs.\ \NSplitRankSeventyFive), although its mean loss is still lower than that of the preset and of plain PSO.
% CHECK-FINAL [C59]: omega = 0.9 vs 0.75: W = L = 0 from either side (\NSplitNinetyVsSeventyFive), case-mean p = \NSplitCmNinetyP >= 0.05, higher mean loss and worse rank than 0.75 but lower mean loss than omega = 1 and 0.5; CHECK-EXTRA [X36]: 0.9 beats 0.75 at no level ("does not help")
In short, of the \NSplitNSettings{} splits, \NSplitBestSetting{} has the lowest mean wake loss and average rank and $\omega=0.9$ does not beat it, so a split above the preset $\omega=0.5$ used here is a candidate to confirm on the target problem. Chosen on \NSplitCases{} cases only, it is not used for the main comparison, which may therefore slightly understate PSO-VNS.
% CHECK-FINAL [C47]: 0.75 best of the \NSplitNSettings settings (lowest mean loss and rank); [C59]: omega = 0.9 vs 0.75 W = L = 0, case-mean p >= 0.05 ("does not beat"); CHECK-EXTRA [X36]: 0.9 beats 0.75 at no level. SWEVO: the split table is back in the main text (tab:split; logged in optA/sw/moved.md)
